\documentclass[10pt,a4paper]{amsart}
\usepackage[margin=19mm]{geometry}
\usepackage{amsmath,amssymb,mathtools}
\usepackage[scr=rsfs,cal=euler]{mathalfa}
\usepackage{xfrac}
\usepackage[unicode,hidelinks,bookmarks=false]{hyperref}
\newcommand{\ie}{i.e.\ }
\newcommand{\cf}{cf.\ }
\newcommand{\resp}{respectively\ }
\newcommand{\Subsection}{\subsection}
\providecommand{\nobreakdash}{\nobreak\hbox{-}}
\setlength{\emergencystretch}{2em}
\multlinegap0pt
\usepackage{float}
\usepackage{xcolor}
\usepackage{amssymb}
\usepackage{tikz}
\usepackage{enumerate}
\usepackage{bm}
\usepackage{mathtools}
\usepackage{xfrac}
\newtheorem{lemma}{Lemma}
\newcommand{\R}{\mathbb{R}}
\newcommand{\calM}{\mathcal{M}}
\newcommand{\diff}{\operatorname{Diff}} %diffeomorphism group, can be decorated any way wished
\DeclareMathOperator{\id}{id}
\title{Khovanov monodromy groups via motions}
\author{Gabriel Corrigan, Livio Ferretti, Isacco Nonino, Susanna Terron}
\date{Source: arXiv:2609.03932v1 (CC BY 4.0)}
\begin{document}
\maketitle

\noindent\emph{Source and licence.} Khovanov monodromy groups via motions. Version arXiv:2609.03932v1 (CC BY 4.0), distributed by arXiv under the Creative Commons Attribution 4.0 International licence, http://creativecommons.org/licenses/by/4.0/, as recorded in the arXiv licence metadata for this exact version and shown on the abstract page https://arxiv.org/abs/2609.03932v1. Full licence notice: \texttt{licenses/arxiv-2609.03932-CC-BY-4.0.txt}.\par\smallskip

\noindent\textit{Editorial scope.} Counted unit: the article's lemma lem:homotopic\_ambient\_isotopy (source0.tex lines 681--683). The article's complete original proof of it is delivered with it (source0.tex lines 685--707, one \texttt{proof} environment of 23 lines). No mathematics is written, repaired or re-derived: the statement and the proof are the article's own lines, in the article's own notation, and no retained line is substituted or re-flowed.  No further source-local context is needed: every cross-reference of every retained line resolves inside the two delivered spans.  Nothing was omitted from any retained span, so the package carries no omission record.  No retained line contains a citation, so the wrapper carries no bibliography and no bibliography entry.  No retained line contains a figure environment, an image-inclusion command or drawing code, so no image is delivered and the package declares no asset dependency.  Editorial formatting only, in addition: an AMS \texttt{amsart} preamble that restates the article's own declarations and macros used by the retained lines, namely the environments \texttt{lemma} on separate counters, together with the standard packages those lines need.  The retained environments print in this document as Lemma 1 in order of appearance, whereas the article prints the same items with the numbering of its own full document.  The complete native source of the article as distributed in the arXiv e-print for this exact version is bundled unchanged as \texttt{sources/209326/source0.tex}.
\begin{lemma}\label{lem:homotopic_ambient_isotopy}
    To every motion~$\gamma \in\calM(L)$ we can associate a cobordism~$\Sigma_\gamma$, well-defined up to ambient isotopy relative to the boundary.
\end{lemma}

\begin{proof}
    Given a path~$\gamma = \{\gamma_s\}_{s\in[0,1]}$ in~$\diff_c(\R^3)$ representing a motion in~$\calM(L)$, let us denote by~$\Sigma_\gamma$ the trace~$(\gamma_s(L),s)$ in~$\R^3 \times [0,1]$. Then~$\Sigma_\gamma$ defines a cobordism from~$L$ to~$\gamma_1(L)=L$.
    If~$\gamma$ and~$\kappa$ are two representatives of the same element of~$\calM(L)$, then there exists a relative homotopy~$H$  between them, i.e.~\begin{align*}
        H \colon [0,1]_t \times [0,1]_s & \to \diff_c (\mathbb{R}^3)\\
         (t,s)  & \mapsto H_t(s)
    \end{align*}
    satisfying
    \begin{align*}
        H_0(s)&=\gamma_s\\
        H_1(s)&=\kappa_s\\
        H_t(0)&=\id_{\mathbb{R}^3}\\
        H_t(1)&\in \diff_c(\mathbb{R}^3,L)
    \end{align*}
    for all~$s, t\in[0,1]$.
    
    Consider the map~$$H' \colon L \times  [0,1]_t\times [0,1]_s \to \mathbb{R}^3 \times [0,1]_s$$
    given by 
    ~$$ H'_t(L,s) = (H_t(s)(L),s).$$
    The map~$H'$ is an isotopy between the embeddings of~$L \times [0,1]_s \hookrightarrow \mathbb{R}^3 \times [0,1]$ given by~$H'_0$ and~$H'_1$, which correspond to the cobordisms~$\Sigma_\gamma=(\gamma_s(L),s)$ and~$\Sigma_{\kappa}=(\kappa_s(L),s)$ traced by the two paths.
    Note that the boundary of the cobordisms is fixed throughout the whole isotopy:~$$H'_t(L,0)= (H_t(0)(L),0)=(L,0),$$
    and~$$ H'_t(L,1)=(H_t(1)(L),1)=(L,1).$$
    Using the isotopy extension theorem, we obtain a map~$$ \overline{H} \colon \mathbb{R}^3 \times [0,1]_t\times [0,1]_s \to \mathbb{R}^3 \times [0,1]_s$$ which is an ambient isotopy between the two surfaces and fixes the boundary of the cobordisms.     
\end{proof}

\end{document}
